\section{Linear estimates on finite orbit segments}\label{sec:linear}

We now turn the diagonal control of
\Cref{prop:good-frames} into estimates for the full derivative.
There are three steps.  A finite-interval change of basis separates
groups of diagonal rates, short derivatives identify nearly the same
spaces at nearby base points, and a product-expansion estimate controls
small perturbations relative to the possibly small exterior volumes.
These statements will supply the linear data for switching orbit
segments.

\subsection{Triangular products and separation of rates}

For invertible matrices $A_t$ indexed by $-n\leq t<n$, write
\[
 A(t,s)=A_{t-1}\cdots A_s\quad(t>s),\qquad
 A(s,s)=\Id,\qquad A(s,t)=A(t,s)^{-1}.
\]
All vector norms and exterior-power norms in the next proposition are
Euclidean operator norms.  If $E$ is a $q$-dimensional subspace and
$A$ is injective, let
\[
 J_q(A|E)=\|(\mathord{\bigwedge}^q A)\xi\|,
\]
where $\xi$ is either unit simple $q$-vector spanning $E$.
This is the intrinsic $q$-dimensional volume expansion.  We set
$J_0=1$, and the zeroth exterior power is the identity on $\R$.

\begin{proposition}[Finite-interval splitting]\label{prop:finite-splitting}
Fix $d\geq1$ and $A_*>1$.  There is a constant $C$, depending only on
these two quantities, with the following properties.
Let $n\geq2$, $b>0$, and $\ln(2n+1)\leq bn$.
Suppose that $A_t$, $-n\leq t<n$, are upper triangular with positive
diagonal and
\[
 \|A_t\|,\ \|A_t^{-1}\|\leq A_*,
 \qquad
 \left|\sum_{t=-n}^{u-1}\ln(A_t)_{jj}
                  -(u+n)\lambda_j\right|\leq bn
\]
for $-n\leq u\leq n$ and $1\leq j\leq d$, where
$\lambda_j=(2n)^{-1}\sum_{t=-n}^{n-1}\ln(A_t)_{jj}$.

First, order the rates as
$\lambda_{(1)}\geq\cdots\geq\lambda_{(d)}$.
For every $0\leq q\leq d$ and every $-n\leq s\leq t\leq n$,
\begin{equation}\label{eq:exterior-rate-comparison}
 \left|\ln\|\mathord{\bigwedge}^q A(t,s)\|
       -(t-s)\sum_{j=1}^q\lambda_{(j)}\right|\leq Cbn.
\end{equation}
The analogous estimate for $A(s,t)$ uses the $q$ largest members of
$\{-\lambda_1,\ldots,-\lambda_d\}$.

Next, suppose that real $\gamma$ and $g>0$ separate all the rates:
each $\lambda_j$ is at most $\gamma-g$ or at least $\gamma+g$.
Let $F$ be the set of indices in the second group, let
$p=|F|$, and put
\[
 \Lambda=\sum_{j\in F}\lambda_j,\qquad
 \Omega=-\sum_{j\notin F}\lambda_j.
\]
There are decompositions
$\R^d=E_t^s\oplus E_t^f$, $-n\leq t\leq n$, invariant under each
$A_t$, with $\dim E_t^f=p$, whose two projections have norm at most
$e^{Cbn}$.  For $s\leq t$ and $l=t-s$ they satisfy
\begin{align}
 \|A(t,s)|E_s^s\|&\leq e^{(\gamma-g)l+Cbn},
 &
 \|A(s,t)|E_t^f\|&\leq e^{-(\gamma+g)l+Cbn},
 \label{eq:split-propagation}\\
 \|\mathord{\bigwedge}^p A(t,s)\|&\leq e^{\Lambda l+Cbn},
 &
 e^{\Lambda l-Cbn}\leq
 J_p(A(t,s)|E_s^f)&\leq e^{\Lambda l+Cbn},
 \label{eq:fast-volume}\\
 \|\mathord{\bigwedge}^{d-p} A(s,t)\|&\leq e^{\Omega l+Cbn},
 &
 e^{\Omega l-Cbn}\leq
 J_{d-p}(A(s,t)|E_t^s)&\leq e^{\Omega l+Cbn}.
 \label{eq:slow-backward-volume}
\end{align}
If $0<p<d$ and at least one of $v_1,\ldots,v_p$ belongs to $E_s^s$,
then
\begin{equation}\label{eq:mixed-exterior}
 \|A(t,s)v_1\wedge\cdots\wedge A(t,s)v_p\|
 \leq e^{(\Lambda-2g)l+Cbn}\prod_{j=1}^p\|v_j\|.
\end{equation}
The backward analogue has exterior degree $d-p$, exponent
$(\Omega-2g)l+Cbn$, and at least one factor in $E_t^f$.
When a group is empty its space is $\{0\}$, all statements on that
space are vacuous, and the exterior-degree-zero conventions apply.
\end{proposition}

\begin{proof}
\emph{Diagonal normalization.}
Write
\[
 e_j(u)=\sum_{t=-n}^{u-1}\bigl(\ln(A_t)_{jj}-\lambda_j\bigr),
 \qquad
 D_u=\diag(e^{e_1(u)},\ldots,e^{e_d(u)}).
\]
The chord hypothesis gives
$\|D_u\|,\|D_u^{-1}\|\leq e^{bn}$.
The matrices
\[
 \widehat A_t=D_{t+1}^{-1}A_tD_t
\]
are upper triangular with constant diagonal
$(e^{\lambda_1},\ldots,e^{\lambda_d})$ and entries bounded by
$e^{Cbn}$.  Since the rates are averages of logarithms of diagonal
entries,
\begin{equation}\label{eq:bounded-rates}
 |\lambda_j|\leq\ln A_*.
\end{equation}
Their inverse matrices also have entries bounded by $e^{Cbn}$.

We record how triangular products are estimated.  If $m$ is fixed and
upper-triangular $m$-by-$m$ matrices have constant diagonal
$e^{\alpha_1},\ldots,e^{\alpha_m}$, bounded $|\alpha_i|$, and
off-diagonal entries at most $e^{Cbn}$, a term in any product entry
can change its coordinate index strictly at most $m-1$ times.
There are at most a fixed constant times $(2n+1)^{m-1}$ choices
for those times and the intermediate indices.
If the product length is $l$ and a term has $a\leq m-1$ strict
changes, its diagonal factors contribute at most
$e^{\alpha_{\max}(l-a)}$.
The omitted $a$ diagonal factors cost at most a fixed factor
$e^{a|\alpha_{\max}|}$ after $e^{\alpha_{\max}l}$ is factored out.
The off-diagonal factors contribute $e^{C_mbn}$.
The polynomial number of terms is absorbed by
$\ln(2n+1)\leq bn$.  Thus the product norm is at most
$e^{\alpha_{\max}l+C_mbn}$.

Apply this argument to $\mathord{\bigwedge}^q\widehat A_t$ in the
coordinate-wedge basis, ordered lexicographically.
This matrix is upper triangular of fixed dimension $\binom dq$,
and its diagonal entries are
$\exp(\sum_{j\in I}\lambda_j)$ for $q$-element index sets $I$.
The largest such sum is $\sum_{j=1}^q\lambda_{(j)}$.
For the reverse inequality, the diagonal of the product indexed by
a maximizing set $I$ is exactly
$\exp(l\sum_{j\in I}\lambda_j)$; that matrix entry is no larger
than the operator norm.  Converting back multiplies the exterior
norms only by the two endpoint gauge bounds.
This proves \eqref{eq:exterior-rate-comparison}.
Applying the same argument to inverse normalized steps in reverse
time proves its backward version.  The cases $q=0,d$ follow as well:
the former is the identity, and the latter is the determinant line.

\emph{Removing entries between the groups.}
We next conjugate $\widehat A_t$ by time-dependent upper
unitriangular matrices to make the two coordinate index groups
invariant.  Process the superdiagonals in increasing order of their
distance $j-i$ from the diagonal.
At a position $i<j$ belonging to different groups, suppose the
current entry is $h_t$.  Use an elementary change
$Q_t=\Id+u_tE_{ij}$, where $E_{ij}$ has its only nonzero entry,
equal to one, in position $(i,j)$.
In $Q_{t+1}^{-1}\widehat A_tQ_t$ the new entry at this position is
\[
 h_t+e^{\lambda_i}u_t-e^{\lambda_j}u_{t+1}.
\]
To set it to zero, solve
\begin{equation}\label{eq:shear-recurrence}
 u_{t+1}=q u_t+f_t,\qquad
 q=e^{\lambda_i-\lambda_j},\qquad f_t=e^{-\lambda_j}h_t.
\end{equation}
If $q<1$, prescribe $u_{-n}=0$ and solve forward:
\[
 u_t=\sum_{a=-n}^{t-1}q^{\,t-1-a}f_a.
\]
If $q>1$, prescribe $u_n=0$ and solve backward:
\[
 u_t=-\sum_{a=t}^{n-1}q^{\,t-1-a}f_a.
\]
In each formula the coefficients have modulus at most one, so
$|u_t|\leq2n\max_a|f_a|$.
Different groups cannot have $q=1$.
This bound avoids dividing by the possibly small separation of rates.

Right multiplication by $Q_t$ changes column $j$ by a multiple of
column $i$, and left multiplication by $Q_{t+1}^{-1}$ changes row
$i$ by a multiple of row $j$.
Upper triangularity shows that every affected off-diagonal position
has distance at least $j-i$; at that distance only $(i,j)$ changes.
The possible quadratic cross term is zero because
$E_{ij}\widehat A_tE_{ij}=(\widehat A_t)_{ji}E_{ij}=0$.
Thus a later operation cannot restore an entry on an earlier
superdiagonal.

There are at most $d(d-1)/2$ operations.
At each operation the bound above, the current entry bounds, and
\eqref{eq:bounded-rates} bound the shear and its inverse by
$e^{Cbn}$.  Multiplying or adding a fixed number of such quantities
only changes $C$; the factor $2n$ is absorbed by $\ln(2n+1)\leq bn$.
Induction over this fixed number of operations therefore gives
matrices $C_t$ with
\begin{equation}\label{eq:finite-gauge}
 \|C_t\|,\ \|C_t^{-1}\|\leq e^{Cbn},
 \qquad
 B_t=C_{t+1}^{-1}A_tC_t,
\end{equation}
where $B_t$ is upper triangular, has diagonal $e^{\lambda_j}$,
and has no entries between the two groups.  Its entries and those of
$B_t^{-1}$ are bounded by $e^{Cbn}$.
All constants are independent of the length of the interval.

Let $\R^F$ be the coordinate span with indices in $F$ and
$\R^{F^c}$ its coordinate complement.  Set
\[
 E_t^f=C_t\R^F,\qquad E_t^s=C_t\R^{F^c}.
\]
These spaces are invariant under the original matrices.
Their projections are conjugates by $C_t$ of coordinate projections,
and hence have norm at most $e^{Cbn}$ after changing $C$.
The two groups may be interlaced in the original coordinate order.
Each induced block is nevertheless upper triangular in its inherited
order; a fixed permutation makes the two blocks consecutive if desired.
The triangular-product estimate within these blocks proves
\eqref{eq:split-propagation}.

\emph{Exterior volumes and mixed factors.}
The fast rates are exactly the $p$ largest rates.  The general exterior
bound already proves the first inequality of \eqref{eq:fast-volume}.
The normalized fast block has determinant $e^{\Lambda l}$ on an
interval of length $l$.
To account explicitly for its possibly oblique basis, define
\[
 V_F(t)=
 \sqrt{\det\bigl((C_t|_{\R^F})^*(C_t|_{\R^F})\bigr)}.
\]
For $p=0$ this Gram determinant is one.  The endpoint bounds in
\eqref{eq:finite-gauge} imply
$e^{-Cbn}\leq V_F(t)\leq e^{Cbn}$.
The intrinsic volume expansion is exactly
\[
 J_p(A(t,s)|E_s^f)
       =e^{\Lambda l}\frac{V_F(t)}{V_F(s)}.
\]
This proves both volume bounds in \eqref{eq:fast-volume}.
Applying the same argument to the inverse slow block gives
\eqref{eq:slow-backward-volume}, including either sign of $\Omega$.

The normalized block decomposition induces an invariant decomposition
of each exterior power according to the number of fast and slow
factors.  In degree $p$, every sector with a slow factor has diagonal
rate sum at most $\Lambda-2g$: at least one fast rate has been
replaced by a slow rate.
The triangular-product estimate therefore bounds this entire sector
by $e^{(\Lambda-2g)l+Cbn}$.
If one original vector is slow, split each of the other vectors into
its fast and slow parts.  There are only a fixed number of terms, each
still having a slow factor.  The projection norms and endpoint gauges
are absorbed into $e^{Cbn}$, proving \eqref{eq:mixed-exterior}.
Reversing time and exchanging the groups gives the backward statement.
The empty-group cases require none of these mixed-factor assertions
and follow from the stated conventions.
\end{proof}

The estimates concern every subinterval of one good finite trajectory.
The change of basis is made once on that trajectory; only its values
at the two ends of a subinterval enter a product estimate.
In particular, a gauge loss is not multiplied once for each time step.

For a good frame in \Cref{prop:good-frames}, the matrices of the
derivative in the moving orthonormal frames satisfy the hypotheses
above for all sufficiently large $k$.
The proposition therefore gives tangent spaces
$E_t^s,E_t^f\subset T_{S^t x}M$ and the corresponding derivative
estimates.  We use this notation also after expressing the spaces in
chart coordinates.  Coordinate maps from a fixed uniformly regular
finite atlas conjugate every product only at its endpoints, so their
bounded norms and inverse norms change the constants in the estimates,
but not their form.

\subsection{Spaces that can be chosen uniformly in a spatial bin}

We need spaces at nearby starting points to agree much more accurately
than their possible small angle.  The comparison uses short derivative
products and the parameters in \eqref{eq:smooth-scales}.
Fix a finite family of coordinate charts with smaller relatively
compact subcharts covering $M$.
Their coordinate maps, inverse maps, and the local expressions of
$S$ and $S^{-1}$ have uniform bounds through order two on the relevant
compact sets.  Each spatial bin below is contained in one of the
smaller subcharts, has coordinate diameter at most $C h_0$, and
uses that specified chart to identify its tangent spaces with $\R^d$.
At a boundary of $M$ use nested half-box charts.

For equal-dimensional subspaces $E,F\subset\R^d$, measure their
distance by $\|\Pi_E-\Pi_F\|$, where $\Pi_E,\Pi_F$ are Euclidean
orthogonal projections.  This distance concerns the subspaces and
does not require a choice of bases in them.

\begin{lemma}[Common spaces in a bin]\label{lem:bin-spaces}
For all sufficiently large $k$, each spatial bin $U$ meeting
$\mathcal G_k$ has a number
\[
 \gamma_U\in[g,(8d+1)g]
\]
such that every good witness frame over $U$ has all its rates at most
$\gamma_U-g$ or at least $\gamma_U+g$, with the same number $p_U$
in the second group.
Apply \Cref{prop:finite-splitting} using this threshold.
There is $c>0$, independent of the bin and the witnesses, such that:
\begin{enumerate}[label=\textup{(\roman*)}]
 \item The spaces $E_0^s$ from any two good witnesses over $U$ are
 within $e^{-c gL}$ of one another in its coordinate chart.
 The same holds for the spaces $E_0^f$.
 \item For good paths starting in $U$ whose time-$n$ endpoints lie
 in a common endpoint bin $V$, their spaces $E_n^f$ are within
 $e^{-c gL}$ of one another in the chart of $V$.
 For paths whose time-$(-n)$ endpoints lie in a common bin,
 the same assertion holds for $E_{-n}^s$.
\end{enumerate}
In particular, fix one representative space for each nonempty
combination of bins in \textup{(ii)}.  Orthogonal projection onto it,
restricted to any corresponding actual endpoint space, is an
isomorphism with inverse norm at most $2$.
If $p_U=0$ or $p_U=d$, the relevant zero or full spaces agree exactly.
\end{lemma}

\begin{proof}
\emph{Short products in common charts.}
Choose one reference point in a starting bin.
Uniform Lipschitz bounds for $S$ and $S^{-1}$ give, for another point
in that bin,
\[
 d_M(S^t x,S^t y)\leq C A_0^{|t|}h_0
 \qquad(|t|\leq L)
\]
with a fixed $A_0>1$.
Since $L=o(Bn)$, this tends to zero uniformly.
At each time on either short reference orbit choose a chart whose
smaller subchart contains the reference point.
All nearby orbits then remain in the same larger chart at that time.
Use the prescribed bin chart at the initial time.
The one-step coordinate derivatives and their inverses are bounded by
a fixed constant, and their differences are at most
$C A_0^{|t|}h_0$.
The telescoping product identity consequently gives
\begin{equation}\label{eq:short-product-comparison}
 \|P_x^\pm-P_y^\pm\|\leq e^{-Bn+CL}=:\zeta_k,\qquad
 \sigma_{\min}(P_x^\pm)\geq e^{-CL}=:m_k.
\end{equation}
Here $P_x^+$ and $P_x^-$ denote $DS^L_x$ and $DS^{-L}_x$,
respectively, in these common initial and terminal charts.
For example, summing the $L$ terms in the product difference bounds
it by $C L A_1^{2L}h_0$, which has the stated form.
The inverse product bound gives the minimum stretch.
On a manifold with boundary the same proof uses the nested half-box
charts and uniform one-sided smooth bounds; equivalently the finitely
many local expressions may be extended to their surrounding boxes.

The singular values are Lipschitz with respect to operator norm.
Since $\zeta_k/m_k\to0$, \eqref{eq:short-product-comparison} implies,
after increasing the constants if necessary,
\begin{equation}\label{eq:short-log-singular}
 |\ln\sigma_j(P_x^+)-\ln\sigma_j(P_y^+)|
       \leq 2\zeta_k/m_k .
\end{equation}
For completeness, the singular-value Lipschitz bound follows from
the min--max characterization, because $\|Pv\|$ and $\|Qv\|$
differ by at most $\|P-Q\|$ on unit vectors.
The logarithmic inequality then follows from the lower bound $m_k$.

\emph{A threshold for every witness.}
Let $\lambda_{(j)}(e)$ be the sorted rates of a good witness $e$.
Subtracting \eqref{eq:exterior-rate-comparison} for exterior degrees
$j$ and $j-1$ gives
\begin{equation}\label{eq:singular-rate}
 \left|L^{-1}\ln\sigma_j(P_{\pi(e)}^+)
                     -\lambda_{(j)}(e)\right|
       \leq Cbn/L.
\end{equation}
This argument uses the no-gap part of
\Cref{prop:finite-splitting}; a threshold has not yet been chosen.
Fix one good witness $e_*$ in $U$.
Equations \eqref{eq:short-log-singular} and \eqref{eq:singular-rate}
show that every good witness in $U$, including every other witness
over the same point, satisfies
\[
 \max_j|\lambda_{(j)}(e)-\lambda_{(j)}(e_*)|
 \leq Cbn/L+2\zeta_k/(m_kL)=o(g).
\]
Consider the $d+1$ candidate thresholds $(1+8j)g$, $0\leq j\leq d$.
A real number lies within distance $3g$ of at most one candidate.
The $d$ representative rates therefore leave a candidate whose
distance from all of them is greater than $3g$.
For large $k$, all the other sorted rate lists differ by less than
$g$, so this candidate separates every list with a margin greater
than $2g$.  It has the weaker required margin $g$, and the number of
rates above it is constant.  Fix this candidate as $\gamma_U$.

\emph{Comparison of subspaces using one short map.}
Assume first that $0<p_U<d$ and put $q=d-p_U$.
Use the single forward short map $P_*=P_{\pi(e_*)}^+$.
Let $W$ be its bottom $q$ right-singular subspace.
Its top $p_U$ singular values are at least
$\exp((\gamma_U+g)L-Cbn)$ by \eqref{eq:singular-rate}.
If $v$ is a unit vector in the slow space of any good witness $e$,
then \eqref{eq:split-propagation} and
\eqref{eq:short-product-comparison} give
\[
 \|P_*v\|\leq e^{(\gamma_U-g)L+Cbn}+\zeta_k.
\]
Orthogonality of the images of right-singular vectors implies
\[
 \|\Pi_{W^\perp}v\|
 \leq e^{-2gL+Cbn}
       +\zeta_k e^{-(\gamma_U+g)L+Cbn}
 \leq e^{-2gL+Cbn}+e^{-Bn+CL}.
\]
In the last bound, constants absorb
$|\gamma_U|=O_d(g)$ and the other terms bounded by a constant times
$L$.  This is at most $e^{-c_1gL}$ for some $c_1>0$ and large $k$.

To pass from this directed estimate to a subspace estimate, observe
that if $\|\Pi_{W^\perp}|E\|\leq\epsilon<1$ and
$\dim E=\dim W$, then $\Pi_W|E$ is an isomorphism.
The space $E$ is a graph over $W$ of norm at most
$\epsilon/\sqrt{1-\epsilon^2}$.
The corresponding orthogonal projectors differ by at most a fixed
multiple of $\epsilon$ when $\epsilon\leq1/2$.
Thus every slow space is within $C e^{-c_1gL}$ of this same $W$;
the triangle inequality compares any two slow spaces.
No continuity assertion about singular vectors or chosen good frames
has been used.

For fast spaces use the single backward short map
$P_{\pi(e_*)}^-$.
Its bottom $p_U$ right-singular space is the comparison space.
The top $q$ singular values are at least
$\exp(-(\gamma_U-g)L-Cbn)$, whereas a unit fast vector has backward
stretch at most $\exp(-(\gamma_U+g)L+Cbn)$.
The identical calculation gives the error bound
\[
 e^{-2gL+Cbn}
       +\zeta_k e^{(\gamma_U-g)L+Cbn}
 \leq e^{-2gL+Cbn}+e^{-Bn+CL}.
\]
This proves \textup{(i)} after reducing $c$ to absorb fixed factors.
In particular, the error is $o(e^{-C_0bn})$ for every fixed $C_0$,
by $bn=o(gL)$.

\emph{Endpoint comparisons.}
Keep the threshold of $U$.
For paths starting in $U$ and ending at time $n$ in a common bin $V$,
choose one representative path for this pair of bins.
Apply the preceding short-product comparison from their time-$n$
points backward for length $L$.
Their initial closeness for this comparison is supplied by $V$,
not by propagating the starting-bin closeness for $n$ steps.
All intervals $[n-L,n]$ lie inside the good interval, and all paths
use the same group dimension and threshold.
In particular, \eqref{eq:exterior-rate-comparison} applies on each
$[n-L,n]$ with the same full-segment rates, so the singular-value
comparison and the backward argument just proved compare their fast
spaces at $n$. At $-n$, the identical all-subinterval estimate for
forward short maps on $[-n,-n+L]$ gives the slow-space assertion.

Finally, if $E$ and its chosen representative $F$ have equal dimension
and $\|\Pi_E-\Pi_F\|\leq\epsilon$, then for $v\in E$,
\[
 \|\Pi_F v\|\geq\sqrt{1-\epsilon^2}\,\|v\|.
\]
The restriction is onto $F$ by equality of dimensions, and its inverse
has norm at most $(1-\epsilon^2)^{-1/2}\leq2$ for large $k$.
If one group is empty, the associated spaces are identically zero or
the full coordinate space, so all comparisons and projection
conclusions follow directly.
\end{proof}

For later use, choose one good witness in each nonempty starting bin
and orthonormal bases within each of its two spaces at time zero.
The matrix having these two bases as columns has norm bounded by a
dimension-only constant and inverse norm at most $e^{Cbn}$, by the
projection bounds of \Cref{prop:finite-splitting}.
It gives fixed product coordinates for that bin.
Likewise \Cref{lem:bin-spaces} supplies fixed endpoint projections
for each relevant pair of bins.
All these choices are finite at a fixed scale.  Their constancy within
the indicated bins, rather than any continuity of choices between
bins, is what the switching construction will use.

\subsection{Relative stability of exterior products}

The bounds on every unperturbed subinterval also control a perturbed
long product.  This must be a relative estimate: a leading volume can
be exponentially small, so continuity with an absolute error alone is
insufficient.

\begin{lemma}[Exterior-product perturbation]\label{lem:relative-exterior}
Fix $d\geq1$ and $K\geq1$.  Let $0\leq p\leq d$, let $m\geq1$,
and let $A_0,\ldots,A_{m-1}$ and
$\widehat A_0,\ldots,\widehat A_{m-1}$ be real $d$-by-$d$ matrices
such that
\[
 \|A_t\|\leq K,\qquad
 \|\widehat A_t-A_t\|\leq\epsilon\leq1.
\]
Use the same product notation as above for nonnegative indices.
Suppose $a\geq0$, $\sigma\in\R$, and
\[
 \|\mathord{\bigwedge}^p A(t,s)\|
       \leq e^{\sigma(t-s)+a}
 \qquad(0\leq s\leq t\leq m).
\]
There is $C=C(d,K)$ such that
\begin{equation}\label{eq:exterior-perturbation}
 \left\|\mathord{\bigwedge}^p\widehat A(m,0)
             -\mathord{\bigwedge}^p A(m,0)\right\|
 \leq e^{\sigma m+a}
       \left[(1+C\epsilon e^{a+|\sigma|})^m-1\right].
\end{equation}
For $p=0$ the left side is zero.

In particular, use the scales \eqref{eq:smooth-scales}, take
$m\leq2n$, $a\leq C_1bn$, $|\sigma|\leq C_2$, and
$\epsilon\leq C_3e^{-3Bn/2}$, with fixed constants.
Then for every fixed $C_0>0$ the left side of
\eqref{eq:exterior-perturbation} is
\begin{equation}\label{eq:relative-exterior-small}
 o\!\left(e^{\sigma m-C_0bn}\right),
\end{equation}
uniformly for all matrices satisfying these hypotheses.
\end{lemma}

\begin{proof}
The exterior power is multiplicative.  Multilinearity in its $p$
columns, together with the one-step bound $K+1$ for the perturbed
matrices, gives
\[
 E_t:=\mathord{\bigwedge}^p\widehat A_t
             -\mathord{\bigwedge}^p A_t,
 \qquad \|E_t\|\leq C\epsilon.
\]
Expand the product using either the unperturbed exterior matrix or
$E_t$ at each time.  A term with $q\geq1$ insertions has $q+1$
unperturbed intervals whose total length is $m-q$.
Allowing empty intervals, the hypothesis bounds that term by
\[
 e^{\sigma(m-q)+(q+1)a}(C\epsilon)^q
 =e^{\sigma m+a}(C\epsilon e^{a-\sigma})^q.
\]
There are $\binom mq$ choices of the inserted positions.
Sum over $q=1,\ldots,m$ and use $-\sigma\leq|\sigma|$ to prove
\eqref{eq:exterior-perturbation}.  This calculation explains why
negative $\sigma$, as may occur for backward exterior volumes, is
harmless: each removed time step contributes only the displayed
factor $e^{-\sigma}$.

For the final assertion put
$u=C\epsilon e^{a+|\sigma|}$ and use
$(1+u)^m-1\leq mu e^{mu}$.
After division by $e^{\sigma m-C_0bn}$ the bound is at most
\[
 C m\epsilon e^{(2C_1+C_0)bn+C_2}
       \exp\!\bigl(Cm\epsilon e^{C_1bn+C_2}\bigr).
\]
The logarithm of its decisive first factor is at most
$-3Bn/2+O(bn)+\ln(2n)$, which tends to $-\infty$ by
\eqref{eq:smooth-hierarchy}.  The quantity in the second exponential
tends to zero for the same reason.
All constants are fixed, proving the uniform little-oh assertion.
For exterior degree zero the two products are both the identity, so
the result holds without an expansion.
\end{proof}

The lemma applies to forward degree-$p$ products with
$\sigma=\Lambda$ and to backward degree-$(d-p)$ products with
$\sigma=\Omega$ from \Cref{prop:finite-splitting}.
Together with \Cref{lem:bin-spaces}, it supplies uniform linear
and relative-volume estimates for orbit segments whose starting
points occupy one small spatial bin.
